\documentclass[12pt,a4paper]{article}

\usepackage[margin=2.54cm]{geometry}
\usepackage[T1]{fontenc}
\usepackage[utf8]{inputenc}
\usepackage{lmodern}
\usepackage{array}
\usepackage{placeins}
\usepackage{longtable}
\usepackage{enumitem}
\usepackage{needspace}
\usepackage[hidelinks]{hyperref}

\setlength{\parindent}{0pt}
\setlength{\parskip}{0.6em}
\setlist[itemize]{nosep,leftmargin=*}
% Consistent visual structure for each prompt.
\newenvironment{promptentry}{%
    \par\Needspace{7\baselineskip}\addvspace{1.2em}
    \hrule
    \nobreak\vspace{0.5em}
}{\par\addvspace{0.5em}}
\newenvironment{aiprompt}{%
    \begin{list}{}{%
        \setlength{\leftmargin}{1em}%
        \setlength{\rightmargin}{1em}%
        \setlength{\topsep}{0.3em}%
        \setlength{\parsep}{0pt}%
    }%
    \item\relax
}{\end{list}}

\AddToHook{cmd/subsection/before}{\Needspace{10\baselineskip}}

\begin{document}

\begin{center}
    {\Large\bfseries Report of Generative AI Use}
\end{center}

\begin{tabular}{@{}l@{\hspace{1.2em}}l@{}}
\textbf{Student name:} & Manon Yiqin Galian Ng\\
\textbf{Student number:} & 23375206\\
\textbf{Module:} & Computer Vision (CSU44053)\\
\textbf{Assignment:} & Pool ball detection and recognition
\end{tabular}

\section{Generative AI Model Used}

Tools and models used for this assignement:
\begin{itemize}
    \item ChatGPT (GPT-6)
    \item OpenAI Codex (GPT-6.1 Sol Medium \& GPT-6.1 Sol High)
\end{itemize}

\textbf{note:} For simplicity, ``ChatGPT'' refers throughout this report to AI assistance
received through either the ChatGPT interface or Codex.

\section{Description of AI Use}

I used ChatGPT as a technical assistant throughout this assignment to discuss
computer vision methods, review design alternatives, translate ideas into
C++/OpenCV code, troubleshoot implementation problems, and improve the organisation
and wording of my report. 

I initially asked ChatGPT to guide my reasoning without providing a complete solution, as I 
wanted to design the system myself. I proposed the main approaches (e.g. separate detection and recognition stages, automated Hough parameter
testing, reference colour histograms, and the use of unique ball identities) and made the final design 
decisions, while ChatGPT helped me question assumptions, refine ideas and implement them.

ChatGPT supplied substantial code examples and, at points, complete or
corrected implementations, which I reviewed, adapted, tested and evaluated against
the assignment images. I performed the results analysis myself and used ChatGPT
to critique that analysis. 

For the report, I planned the structure and selected the information to include. ChatGPT helped 
me put my thoughts into words, improve clarity and develop the report I had in mind.

\section{Main Uses and Prompts}

\subsection{Understanding the Brief}

\begin{promptentry}

\begin{aiprompt}
"I have a computer vision assignment I am working on. It is important that
you do not give me a solution. I want to design the system by myself. Just act as a guide and help me though my own solution. 
This assignment will use OpenCV and C++. 
I will also provide my lecture notes so we can stay in the scope of the module.
It is important to use traditional computer vision techinique and not a neural network.
I will attach the assignment brief."
\end{aiprompt}

\textbf{Impact on project:} Explanation of the brief, noted that there are separate training and eval
images. Noted that it will not propose full solutions and that we should use traditional CV techniques

\end{promptentry}
\FloatBarrier

\begin{promptentry}

\begin{aiprompt}
"It seems like we cannot use the Table1-6 for training. So that means we can only tweek
parameters on the AllBalls image right? But would we be allowed to look at the validation results
and find general patterns that we are missing and then build new experiments around those errors?"
\end{aiprompt}

\textbf{Impact on project:} Noted that yes we cannot use the Table1-6 for training, e.g. parameter tweaking.
Also if we wanted to structure new experiment around the validation result this should be noted.

\end{promptentry}
\FloatBarrier

\subsection{Help for Detection}

\begin{promptentry}

\begin{aiprompt}
"For this assignment let's first work on the ball detection. I think
the best approach is to use a Hough circle transform."
\end{aiprompt}

\textbf{Impact on project:} Explained that the circle Hough transform would be 
good start and propose starting parameters and what each parameters did.

\end{promptentry}
\FloatBarrier

\begin{promptentry}

\begin{aiprompt}
"We have been trying to find the best Hough manually. This is super time consuming.
I think we should be able to automate this process using the test image.
I want to test a large amount of parameters (a few hundred combinations) without manually testing them all.
If we can, we should also evaluate each config in a specific way.
For the next stage we can only recognise ball that have been detected.
So we should focus on getting pretty much 100\% recall. Then focus on precision
to avoid getting too many FPs. We can validate afterwards on Table1-6."
\end{aiprompt}

\textbf{Impact on project:} Confirmed that automation was appropriate and 
proposed a range of parameters to test (200+, I later increased it to 480 in the final
pipeline after more experimenting). Also proposed the validation structure from my criteria.

\end{promptentry}
\FloatBarrier

\begin{promptentry}

\begin{aiprompt}
"We recently learnt about watershed segmentation. Could we use this method to improve
detection on regions where balls are closely clustered together?
Could you help me implement an experiment for this?"
\end{aiprompt}

\textbf{Impact on project:} Helped implement an experiment for watershed segmentation. 
After testing this experiment the performance did not change much. So this was not included in the
final pipeline. This led to another idea to improve detection; using a more permissive Hough in
the ball clusters and the current Hough everywhere else. This experiment managed to detect 1 missing ball
but added a lot more FP and complexity so was also left out of the final pipeline.

\end{promptentry}
\FloatBarrier

\subsection{Help for Recognition}

\begin{promptentry}

\begin{aiprompt}
"For recognition I initially think we could just use HLS maybe?
I don't think this approach will perform amazingly but I think it 
will be good in the report later for comparaison.
we can look at the pixel inside the balls and divided the pixels
into light, dark, coloured and classify the types based on the proportions.
What do you think?"
\end{aiprompt}

\textbf{Impact on project:} Explained that it was a good simple start. 
Gave some starting boundaries I could test.

\end{promptentry}
\FloatBarrier

\begin{promptentry}

\begin{aiprompt}
"Actually! could we try k-means? with k=2 then hopefully we can extract 2 colour clusters.
Then we can examine what these clusters are (light, dark, colour) using some
thresholds. Then we can look at the proportion of these colour regions for each
ball and see if they are stipped, full or the cue ball. 
full should be primariy dark or colour pixels, stripped should have a larger proprotion of light and 
have dark/colour pixels and the cue ball should be all or mostly light.
Then we could use Hue to get the colour of the ball. 
We can also test other k sizes, k=2 just seem like the simplest and most logical."
\end{aiprompt}

\textbf{Impact on project:} Helped implement this k means experiment and the visualisations.
We tested with k = 2,3,4,5. 
With k=2 we were getting optimal performance on the training image so there was not enough evidence to switch to a 
larger k.

\end{promptentry}
\FloatBarrier

\begin{promptentry}

\begin{aiprompt}
"We did a tutorial where we had to detect taffic cones (the orange and white ones). 
We used at some point to see if there was a white part 'sandwhiched' between two orange 
regions. For stripe balls this could be useful maybe, at least when the orientation is good.
Essentially, stripe balls should have 3 main regions with two white and one colour in the middle.
Can you help me write an experiment to test this?"
\end{aiprompt}

\textbf{Impact on project:} Also helped me plan and implement this experiment. 
Some initial tests on the validation image were performing worse. So it was left as an experiment
and not included in the final report.

\end{promptentry}
\FloatBarrier

\begin{promptentry}

\begin{aiprompt}
"maybe since we have the training image we could use histogram backprojection / Bhattacharyya for better 
colour recogntion? We could use the AllBalls image for reference.
This could help capture some variations taht the single HSV rules cannot capture."
\end{aiprompt}
\FloatBarrier

\textbf{Impact on project:} AI's most substantial contribution here was helping me plan and implement
the histogram training based on the colours of the balls in the training image. 
Help a lot with the code implementation.

\end{promptentry}
\FloatBarrier

\begin{promptentry}

\begin{aiprompt}
"I am thinking. One thing that we have not yet taken into account is that balls are unique. 
Thus, instead of having 'hard' decisions for each ball we can have probabilities or scores.
Then we can pick based on the combination that maximises the probabily or scores.
I think there is an algorithm that does this but I do not remeber the name or how it works 
exactly"
\end{aiprompt}

\textbf{Impact on project:} Gave the name of the algorithm I was missing (Hungarian algorithm)
proposed a Gaussian function for the scores. Helped implement the code for the algorithm.

\end{promptentry}
\FloatBarrier

\begin{promptentry}

\begin{aiprompt}
"Yes that makes sense. A Gaussian function makes sense, we could also try a linear one first. 
But probably the smooth decisions that the Gaussian would give a better performace than decisions 
based on a threshold that the linear function would give. We also need to find an appropriate gaussian form. 
What do you think?"
\end{aiprompt}

\textbf{Impact on project:} We implemented and compared linear and Gaussian compatibility 
functions. The Gaussian formulation was the only one retained in the final implementation 
and iterated on. ChatGPT assisted with implementing both approaches.

\end{promptentry}
\FloatBarrier

\begin{promptentry}

\begin{aiprompt}
"Another thing that could be interesting is that, specifically for duplicated cue balls.
We could expand on pixels around the detection and see if there are colour pixels beyond the white
(ofc this would be tough for green balls given the green mat).
This is because a common error, is that sometimes the hough transform detects the smaller inside
white circle of the balls instead of the full ball."
\end{aiprompt}

\textbf{Impact on project:} We tested multiple padding 1-5 pixels. 1 pixel worked the 
best and is what was kept.

\end{promptentry}
\FloatBarrier

\begin{promptentry}

\begin{aiprompt}
"The padding is working well. One last thing I would like to try, 
is to use the fact that all ball have roughly the same size.
So is there are still duplicate white balls then we could use the estimate 
size at a certain dept to determine which ball is the acctuall white cue.
Additionally, to make it simplier we could also use a simple heuristic that
says that the larger of the two should be the correct white ball. Then we don't 
have to compute estimate ball sizes at all depths."
\end{aiprompt}

\textbf{Impact on project:} We tested both the simpler and more complex 
solution. The complex solution involved taking the training image and fitting a line 
relating detected radius to the center row using the 16 balls. We then divided each
candidate radius by the expected radius at its row. 
Both methods produced the same improvements so we kept the simpler radius rule.

\end{promptentry}
\FloatBarrier

\subsection{Help for Evaluation}

\begin{aiprompt}
"maybe in the evaluation we can also evaluate each ball type like how well it evaluates all the red (an also type like stripe and full)
then we can acc see where the system has some difficukties.
help me write a function in the chat do to this"
\end{aiprompt}

\textbf{Impact on project:} Helped implement a function that help me pinpoint better the shortcomings of the system
that was designed.

\subsection{Help for Result Visualisation}

\begin{aiprompt}
"could you help me write a function to created a top view grid of the results of the system.
I think that would be more clear than what I currently have. Which is a top view of the table 
with the detected circle and the label"
\end{aiprompt}

\textbf{Impact on project:} Helped implement a function to draw this results grid with the detected and recognised balls.

\subsection{Help for Report}

\begin{aiprompt}
"in the Ball type recognition section I added some images but I would also like a nice visulisation. 
to support the pictures I added, maybe someting to show to final classification?"
\end{aiprompt}

\textbf{Impact on project:} made the small visualisation in figure 4

\begin{aiprompt}
"I want to illustration the histogram referencing the report. I think we could draw
approximative histogram for a ball pair and then show the averaging that is done.
Also for that section I want a nice visulisation for how this approaches the performance.
we could take 1-3 balls that were previsouly miss classified and how the correction. Makes sure
to show that classification works on scores now and not hard boundaries."
\end{aiprompt}

\textbf{Impact on project:} made the visualisation for figure 5 and 6.

\begin{aiprompt}
"Given the ball by ball performace: [performance output]. Help me create a visulisation for this.
I was thinking a bar-like chart with the annotated vs detected vs correct identity recognition for each
ball"
\end{aiprompt}

\textbf{Impact on project:} made the visualisation for figure 9.

\textbf{Additionally, some repetitive prompts such as:}

\begin{aiprompt}
"Help me rephrase this" 
\end{aiprompt}

\begin{aiprompt}
"Help me make this sound better"
\end{aiprompt}

\subsection{Help for Code/Syntax}

\textbf{Used some repetitive prompts such as:}

\begin{aiprompt}
"Help me debug this"
\end{aiprompt}

\begin{aiprompt}
"Help me restructure this code in a nicer way"
\end{aiprompt}

\begin{aiprompt}
"Help me write a function for ..."
\end{aiprompt}

\begin{aiprompt}
"Help me draft an experiment for ..."
\end{aiprompt}

\begin{aiprompt}
"How can I implement this in OpenCV?"
\end{aiprompt}

\begin{aiprompt}
"I don't like your implementation maybe we can try ..."
\end{aiprompt}

\begin{aiprompt}
"this function can definitly be a lot shorter and more efficient"
\end{aiprompt}

\subsection{Other:}

\textbf{Used some repetitive prompts such as:}

\begin{aiprompt}
"What does this mean?"
\end{aiprompt}

\begin{aiprompt}
"Could you please clarify this?"
\end{aiprompt}

\section{Coverage checklist}

Summary of AI usage in the assignment.

\begin{itemize}
    \item Planning the detection and recognition pipeline.
    \item Generating, modifying, consolidating or formatting C++ code.
    \item Debugging and explaining code or algorithms.
    \item Designing experiments, parameter searches and interactive tools.
    \item generating LaTeX template for the report and AI usage document
    \item Proofread the report and helped re-formulate some section for added clarity
    \item Helped generate some visualisation for the report with gidence and from the given results
\end{itemize}

\end{document}
